%

Wzmocnienie nierówności Eulera: % lang-pl
Un rafforzamento della disuguaglianza di Eulero: % lang-it

\begin{proposition}%
[nierówność Mitrinovica] % lang-pl
[disuguaglianza di Mitrinović] % lang-it
\index{nierówność!Mitrinovica}% % lang-pl
\index{disuguaglianza!di Mitrinović}% % lang-it
    Dany jest trójkąt o bokach $a, b, c$, wpisany w okrąg o promieniu $R$, opisany na okręgu o promieniu $r$. % lang-pl
    Dato un triangolo di lati $a$, $b$, $c$, inscritto in una circonferenza di raggio $R$ e circoscritto a una circonferenza di raggio $r$. % lang-it
    Wtedy % lang-pl
    Sia % lang-it
    % https://www.cut-the-knot.org/m/Geometry/InequalitiesInTriangle.shtml
    \begin{equation}
        6 \sqrt{3} r \le a + b + c \le 3 \sqrt{3} R.
    \end{equation}
\end{proposition}

Guzicki \cite[s. 362, 363]{guzicki_2021} dowodzi tego z nierówności izoperymetrycznej oraz oszacowania % lang-pl
Guzicki \cite[p. 362, 363]{guzicki_2021} dimostra questo risultato a partire dalla disuguaglianza isoperimetrica e dalla stima % lang-it
\begin{equation}
    \sin \alpha + \sin \beta + \sin \gamma \le 3 \sqrt{3} / 2.
\end{equation}

Marcin Pilipczuk wywnioskuje tę nierówność w $\Delta_{08}^{6}$ z nierówności Jensena dla funkcji wypukłej $-\sin x$ na przedziale $[0, \pi]$ i zauważy, że spośród trójkątów wpisanych w okrąg o promieniu $1$ największy obwód ma trójkąt równoboczny (z twierdzenia sinusów obwód wynosi $2 \sin \alpha + 2 \sin \beta + 2 \sin \gamma$). % lang-pl
Marcin Pilipczuk dedurrà questa disuguaglianza in $\Delta_{08}^{6}$ dalla disuguaglianza di Jensen per la funzione convessa $-\sin x$ sull'intervallo $[0, \pi]$ e osserverà che, tra i triangoli inscritti in una circonferenza di raggio $1$, quello di perimetro massimo è il triangolo equilatero (per il teorema dei seni il perimetro è $2 \sin \alpha + 2 \sin \beta + 2 \sin \gamma$). % lang-it

% TODO: https://www.cut-the-knot.org/m/Geometry/InequalitiesInTriangle.shtml

%